\documentclass[11pt]{article}
\usepackage[margin=28mm]{geometry}
\usepackage{amsmath,amssymb,amsthm}
\usepackage[T1]{fontenc}
\usepackage{hyperref}
\hypersetup{colorlinks=true,linkcolor=blue,citecolor=blue,urlcolor=blue}
\newtheorem{theorem}{Theorem}
\newtheorem{lemma}[theorem]{Lemma}
\newcommand{\tr}{\operatorname{tr}}
\newcommand{\Lam}{\Lambda_+}
\title{The largest-root extremizer for hyperbolic cubics}
\author{OpenAI \texttt{gpt-6.1-sol}, reasoning \texttt{ultra}}
\date{28 October 2026}
\begin{document}
\maketitle
\begin{abstract}
We give a self-contained proof of the concentrated-scalar extremizer in the
mixed-characteristic largest-root problem for every homogeneous hyperbolic
polynomial of degree three. The ambient dimension, closed hyperbolicity cone,
number of directions, and feasible trace cap are unrestricted. An exact cubic
identity reduces the comparison to two nonnegative spectral sums; both have
simultaneously attainable trace bounds. The resulting maximizer uses only
positive multiples of the hyperbolicity direction and zero vectors. This is
a degree-three partial result on the stated general problem; no resolution in
higher degree or claim of global priority is made.
\end{abstract}

\section{Problem, checked antecedents, and scope}

Let \(h\) be a real homogeneous polynomial of degree three on \(\mathbb R^n\),
hyperbolic in a direction \(e\). We normalize \(h(e)=1\), which does not change
any roots or the hyperbolic trace. Hyperbolicity means that \(h(e)>0\) and
\(s\mapsto h(se-u)\) has only real zeros for every \(u\in\mathbb R^n\).
Write these zeros as \(\lambda_1(u),\lambda_2(u),\lambda_3(u)\). Define
\[
 \Lam=\{u:\lambda_j(u)\geq0\text{ for }j=1,2,3\},\qquad
 \tr(u)=\sum_{j=1}^3\lambda_j(u).
\]
Thus \(\Lam\) is the closed hyperbolicity cone and \(e\) is in its interior.
For a direction \(u\), let \(D_u h(x)=\left.\frac{d}{ds}h(x+su)\right|_{s=0}\).
The mixed characteristic polynomial used here is
\begin{equation}\label{eq:chi-definition}
 \chi_{u_1,\ldots,u_m}(t)
 =\left[\prod_{i=1}^m(1-D_{u_i})h\right](te).
\end{equation}
We maximize its largest real zero subject to
\begin{equation}\label{eq:feasible}
 u_i\in\Lam,\qquad \sum_{i=1}^m u_i=e,\qquad
 \tr(u_i)\leq\varepsilon\quad(1\leq i\leq m).
\end{equation}

The matrix version of the concentrated-scalar conjecture appears in the
conclusion of Marcus, Spielman, and Srivastava~\cite{MSS}.
Br\"and\'en~\cite[Section 6, Problem 1 and Conjectures 6.1--6.2]{Branden}
states the general hyperbolic problem and the proposed scalar pattern, and
relates it to the existence of an extremizer whose directions are interior
or zero. The cited workshop question~\cite[p.~995]{OWR} asks this latter
question. Br\"and\'en's Proposition 6.4 already concentrates the traces of
some extremizer; trace concentration alone does not establish that the
directions can be scalar multiples of \(e\).

The pinned repository attempt~\cite{Prior} proves an exact calculation only
for \(h(x,y)=xy\). Here we treat every hyperbolic cubic, including cones
without an assumed determinantal representation. The proof below is
independent of any such representation and of Proposition 6.4.

\section{Exact maximum in degree three}

For \(\varepsilon>0\), set
\begin{equation}\label{eq:constants}
 k=\left\lfloor\frac{3}{\varepsilon}\right\rfloor,
 \quad r=3-k\varepsilon,\quad
 Q_* =\frac{k\varepsilon^2+r^2}{3},\quad
 R_* =\frac{k\varepsilon^3+r^3}{27}.
\end{equation}
In particular \(0\leq r<\varepsilon\), and \(Q_*,R_*>0\).

\begin{theorem}\label{thm:main}
Let \(h,e\) be as above. Let \(m\geq1\) be an integer and
\(\varepsilon>0\). The constraint set \eqref{eq:feasible} is nonempty if and
only if \(m\varepsilon\geq3\). When it is nonempty, the maximum largest zero
of \eqref{eq:chi-definition} is \(1+Y_*\), where \(Y_*\) is the unique
positive zero of
\begin{equation}\label{eq:optimal-cubic}
 Y^3-Q_*Y-2R_*=0.
\end{equation}
It is attained by \(k\) directions equal to \(\varepsilon e/3\), one
direction equal to \(re/3\) if \(r>0\), and all remaining directions zero.
Consequently every nonzero direction of this maximizer is in the interior
of \(\Lam\).
\end{theorem}

\begin{lemma}\label{lem:identity}
For any feasible tuple, put
\[
 q(u)=\sum_{1\leq a<b\leq3}\lambda_a(u)\lambda_b(u),\qquad
 Q=\sum_{i=1}^m q(u_i),\qquad R=\sum_{i=1}^m h(u_i).
\]
Then \(Q,R\geq0\), and
\begin{equation}\label{eq:cubic-identity}
 \chi_{u_1,\ldots,u_m}(t)=(t-1)^3-Q(t-1)-2R.
\end{equation}
\end{lemma}
\begin{proof}
Homogeneity and the eigenvalue factorization give the polynomial identity
\begin{equation}\label{eq:direction-expansion}
 h(e+su)=\prod_{j=1}^3(1+s\lambda_j(u))
 =1+s\tr(u)+s^2q(u)+s^3h(u).
\end{equation}
In particular \(\tr(u)=D_u h(e)\) is linear in \(u\),
\(D_u^2h(e)=2q(u)\), and \(D_u^3h=6h(u)\).
The nonnegativity follows from the nonnegative eigenvalues on \(\Lam\).

Write \(D_i=D_{u_i}\). The operators commute and \(\sum_iD_i=D_e\).
Terms of order greater than three vanish in \eqref{eq:chi-definition}.
The order-zero and order-one terms are \(t^3\) and \(-3t^2\).
The identity
\[
 \sum_{i<j}D_iD_j
 =\frac12\left(D_e^2-\sum_iD_i^2\right)
\]
gives the order-two term \((3-Q)t\).
For order three, expansion of the commuting operators yields
\begin{equation}\label{eq:operators}
 6\sum_{i<j<\ell}D_iD_jD_\ell
 =D_e^3-3\sum_iD_i^2D_e+2\sum_iD_i^3.
\end{equation}
The polynomial \(D_i^2h\) is linear, so
\(D_i^2D_e h=D_i^2h(e)=2q(u_i)\). Also \(D_e^3h=6\).
Applying \eqref{eq:operators} to \(h\) shows that the order-three
coefficient before its negative sign is \(1-Q+2R\). Thus
\[
 \chi(t)=t^3-3t^2+(3-Q)t-1+Q-2R,
\]
which is \eqref{eq:cubic-identity}. Empty pair or triple sums cause no
exception when \(m<3\).
\end{proof}

\begin{lemma}\label{lem:bounds}
Every feasible tuple satisfies \(Q\leq Q_*\) and \(R\leq R_*\).
The scalar tuple in Theorem~\ref{thm:main} attains both equalities.
\end{lemma}
\begin{proof}
Let \(\tau_i=\tr(u_i)\). Linearity of trace implies
\(\sum_i\tau_i=\tr(e)=3\); moreover \(0\leq\tau_i\leq\varepsilon\).
For three nonnegative numbers with sum \(\tau_i\), the sum of pairwise
products is at most \(\tau_i^2/3\), by the sum of squared pairwise
differences. Their product is at most \(\tau_i^3/27\), by the
arithmetic--geometric mean inequality. Therefore
\begin{equation}\label{eq:spectral-bounds}
 Q\leq\frac13\sum_i\tau_i^2,\qquad
 R\leq\frac1{27}\sum_i\tau_i^3.
\end{equation}

We give the simultaneous trace maximization explicitly. If two entries
\(a,b\) satisfy \(0<a\leq b<\varepsilon\), transfer
\(\delta=\min\{a,\varepsilon-b\}>0\) from \(a\) to \(b\).
For \(p=2,3\), convexity of \(x^p\) on \([0,\infty)\) shows that
\((a-\delta)^p+(b+\delta)^p\geq a^p+b^p\): along the transfer the
derivative is \(p((b+s)^{p-1}-(a-s)^{p-1})\geq0\).
At least one of the two entries reaches a boundary. Repeating finitely
many times leaves at most one entry strictly between \(0\) and
\(\varepsilon\), without decreasing either sum. The total mass forces
the resulting entries to be \(k\) copies of \(\varepsilon\), the
optional remainder \(r\), and zeros. This proves the two upper bounds.

For \(u=(\tau/3)e\), all three eigenvalues equal \(\tau/3\), so both
inequalities in \eqref{eq:spectral-bounds} are equalities. The stated
scalar tuple has sum \(e\) and attains \(Q_*,R_*\).
It fits in \(m\) positions: \(m\varepsilon\geq3\) gives \(k\leq m\),
and gives \(k+1\leq m\) when \(r>0\). If \(r=0\), no remainder
position is used. If \(\varepsilon>3\), then \(k=0,r=3\), so the tuple
consists of \(e\) and zeros.
\end{proof}

\begin{proof}[Proof of Theorem~\ref{thm:main}]
Necessity of \(m\varepsilon\geq3\) follows by summing the trace caps.
The scalar tuple constructed in Lemma~\ref{lem:bounds} proves sufficiency.

For \(Q,R\geq0\), consider \(f_{Q,R}(Y)=Y^3-QY-2R\).
If \(R>0\), this polynomial is negative for
\(0\leq Y\leq\sqrt Q\) and strictly increasing from a nonpositive
value to infinity on \([\sqrt Q,\infty)\). Hence it has a unique
positive zero, larger than \(\sqrt Q\), and this is its largest real
zero. If \(R=0\), its largest real zero is \(\sqrt Q\), also when
\(Q=0\). Thus in all cases the largest real zero lies on that
nondecreasing branch.

Let \(Y_*\) be the positive zero for \(Q_*,R_*\). Then
\(Y_*>\sqrt{Q_*}\geq\sqrt Q\), and
\[
 f_{Q,R}(Y_*)=(Q_*-Q)Y_*+2(R_*-R)\geq0.
\]
The branch observation implies that the largest real zero of
\(f_{Q,R}\) is at most \(Y_*\). Lemma~\ref{lem:identity} transfers
this comparison to the zeros in \(t\), adding one. The scalar tuple
attains both coefficient bounds and hence the claimed maximum.
Every positive multiple of \(e\) is in the interior of \(\Lam\),
which proves the last assertion.
\end{proof}

\section{Consequences and limitations}

Theorem~\ref{thm:main} establishes the prescribed extremizer for the entire
degree-three class of the original hyperbolic problem, rather than for a
single product polynomial. In particular it answers the interior-or-zero
question affirmatively in that class. The maximum depends on the trace cap,
and on \(m\) only through feasibility. For example, at \(\varepsilon=1\)
and \(m\geq3\), the exact maximum is \(1+Y_*\) with
\(Y_*^3-Y_*-2/9=0\), attained by three copies of \(e/3\).
For \(\varepsilon\geq3\), it is \(3\), attained by the single nonzero
direction \(e\).

No assertion is made that all maximizing tuples must be scalar. No result
for degrees four and higher is proved. In those degrees the differential
expansion contains additional mixed terms; the two simultaneously monotone
coefficients used here do not constitute a proof there. Nor does this
argument establish a majorization comparison of all roots.

The literature check covered the original workshop passage, the stated
conjectures and accompanying arguments in \cite[Section 6]{Branden}, and
the matrix conjecture in \cite[Section 7]{MSS}. These checked passages
do not supply the cubic proof above. This is not an exhaustive priority
search. In particular Samuel Zackrisson's 2017 KTH thesis,
\emph{Coefficients and zeros of mixed characteristic polynomials}
(DiVA record \texttt{diva2:1106332}), was identified but its full text could
not be retrieved in this investigation. No claim about that thesis's
contents, or about the first appearance of the cubic result, is made.
The claim is a complete proof of the specified degree-three case, with
global novelty unverified.

\section*{Provenance and attribution}

The repository target is
\texttt{prb\_465ecdb5ab0947d19b52e95361b1df42}, statement revision
\texttt{stm\_2c00dc6dd9a74da08a9ba7dc55cb24bb}, read at commit
\texttt{c4ef184d1d392f70fed32c5e862d7f24d11e8401}. The previous
contribution is \texttt{con\_6b719153bc2b4375b48fbd2b661087aa}, revision
\texttt{rev\_c4c3320510fc413eb308160bffb7a3cd}.
The mathematical derivation and exposition were generated by OpenAI
\texttt{gpt-6.1-sol}, reasoning \texttt{ultra}, research agent b.
Research agent c, configured with the same model and reasoning setting,
independently rederived the cubic identity and checked the preliminary
proof route; final independent review is recorded separately. Codex based
on GPT-6 coordinated the research. Exact serving builds and inference time
budgets were not exposed. There is no claim of human authorship,
independent human review, or formal proof verification.

No LaTeX engine was available to the author; compilation and rendered
layout inspection were not performed. Static source checks and independent
source-format reviews are recorded separately. Original exposition is
offered under CC BY 4.0; cited sources retain their own rights.

\begin{thebibliography}{9}
\bibitem{Branden}
P. Br\"and\'en,
\emph{Hyperbolic polynomials and the Marcus--Spielman--Srivastava theorem},
arXiv:1412.0245, Section 6.
\url{https://arxiv.org/abs/1412.0245}.

\bibitem{MSS}
A. W. Marcus, D. A. Spielman, and N. Srivastava,
\emph{Interlacing families II: Mixed characteristic polynomials and the
Kadison--Singer problem}, arXiv:1306.3969v4, Section 7.
\url{https://arxiv.org/abs/1306.3969v4}.

\bibitem{OWR}
\emph{Real Algebraic Geometry With A View Toward Systems Control and Free
Positivity}, Oberwolfach Reports, DOI 10.4171/owr/2014/17,
pp. 977--1045; relevant question on p.~995.
\url{https://doi.org/10.4171/owr/2014/17}.

\bibitem{Prior}
We-Solve/research, archived contribution
\texttt{con\_6b719153bc2b4375b48fbd2b661087aa},
\emph{Interior Maximizers of Mixed Characteristic Polynomials},
degree-two product-cone calculation, pinned repository source.
\url{https://github.com/We-Solve/research/blob/c4ef184d1d392f70fed32c5e862d7f24d11e8401/problems/46/prb_465ecdb5ab0947d19b52e95361b1df42/contributions/con_6b719153bc2b4375b48fbd2b661087aa/files/manuscript.tex}.
\end{thebibliography}
\end{document}
